\documentclass[11pt]{article}
\usepackage[a4paper,margin=27mm]{geometry}
\usepackage[T1]{fontenc}
\usepackage{lmodern,microtype}
\usepackage{amsmath,amssymb,amsthm,mathtools}
\usepackage{enumitem,xcolor}
\usepackage[colorlinks=true,linkcolor=blue!55!black,citecolor=blue!55!black,urlcolor=blue!55!black]{hyperref}
\hypersetup{
 pdftitle={Mixed-Jet Rank Floors for Point-Span Higher Osculating Absorption},
 pdfauthor={Lluis Eriksson},
 pdfsubject={Algebraic geometry; osculating spaces; jets; interpolation},
 pdfkeywords={osculating spaces, jet ampleness, fat points, tangent absorption}
}
\newtheorem{theorem}{Theorem}[section]
\newtheorem{proposition}[theorem]{Proposition}
\newtheorem{lemma}[theorem]{Lemma}
\newtheorem{corollary}[theorem]{Corollary}
\theoremstyle{definition}
\newtheorem{definition}[theorem]{Definition}
\theoremstyle{remark}
\newtheorem{remark}[theorem]{Remark}
\theoremstyle{plain}
\newcommand{\PP}{\mathbf P}
\newcommand{\kk}{\mathbf k}
\newcommand{\cO}{\mathcal O}
\newcommand{\Span}{\operatorname{Span}}
\newcommand{\Osc}{\operatorname{Osc}}

\title{Mixed-Jet Rank Floors for Point-Span\\Higher Osculating Absorption}
\author{Lluis Eriksson\\\small Independent researcher}
\date{August 24, 2026\quad (version 0.2)}

\begin{document}
\maketitle

\begin{abstract}
Let $X$ be a smooth projective integral $d$-fold over an algebraically
closed field, let $H$ be very ample, and use the complete embedding defined
by $H^m$.  Fix $1\leq s\leq m$.  A nonempty finite reduced set
$Z\subset X$ is point-span $s$-osculating-absorbing if the span $S_Z$ of
its value evaluations contains the affine osculating space of order $s$ at
every support.  Write
\[
 m+1=q(s+1)+r,\qquad 0\leq r\leq s,
\]
and put
\[
 B_{d,s}(m)=q\binom{d+s}{d}
 +\begin{cases}0,&r=0,\\ \binom{d+r-1}{d},&r>0.\end{cases}
\]
We prove in arbitrary characteristic that
\[
 \dim S_Z,\ |Z|\geq B_{d,s}(m).
\]
The proof is scheme-theoretic.  Products of point separators isolate
arbitrary mixed jets of total weight at most $m+1$, and discrete convexity
selects full order-$s$ blocks plus one residual block.  The resulting
heterogeneous rank certificate also applies to arbitrary jet-ample
polarizations.  For $s=1$ the formula recovers the earlier growing first-jet
bound.  Equality holds in every characteristic on rational normal curves,
for every $s$, and on $\PP^d$ in the top-order regime $s=m$.  No implication
from tangent absorption to the higher-order hypothesis is used or claimed.
No classification or exhaustive priority claim is made.
\end{abstract}

\section{Statement and scope}

Let $\kk$ be algebraically closed, let $X$ be smooth, projective, integral,
and of dimension $d\geq1$, and let $H$ be very ample.  Write
\[
 \varphi_m:X\longrightarrow\PP(H^0(X,H^m)^*)
\]
for the complete embedding.  For a nonempty finite reduced set $Z\subset X$,
set
\[
 S_Z=\operatorname{Im}\!\left(
 H^0(Z,H^m|_Z)^*\longrightarrow H^0(X,H^m)^*\right).
\]
Equivalently, $S_Z$ is the span of the evaluation lines at the supports;
this formulation makes no choice of fiber trivializations.
For $p\in X$ and $a\geq0$, put
\[
 (a+1)p=\operatorname{Spec}(\cO_{X,p}/\mathfrak m_p^{a+1}),\qquad
 \lambda_d(a)=\binom{d+a}{d},
\]
and set $\lambda_d(-1)=0$.

The order-$a$ principal-parts evaluation is
\[
 j^a_p:H^0(X,H^m)\longrightarrow
 H^0\!\left((a+1)p,H^m|_{(a+1)p}\right).
\]
We define the affine order-$a$ osculating space intrinsically by
\[
 \widehat\Osc^a_p(H^m)=\operatorname{Im}((j^a_p)^*)
 \subset H^0(X,H^m)^*.
\]
This convention uses no ordinary derivatives or factorials.

\medskip

\begin{definition}
Fix $1\leq s\leq m$.  The set $Z$ is
\emph{point-span $s$-osculating-absorbing} for $H^m$ if
\[
 \widehat\Osc^s_p(H^m)\subseteq S_Z\qquad(p\in Z).
\]
\end{definition}

Write uniquely
\[
 m+1=q(s+1)+r,\qquad q\geq1,\quad 0\leq r\leq s,
\]
and define
\begin{equation}\label{eq:B}
 B_{d,s}(m)=q\lambda_d(s)+\lambda_d(r-1).
\end{equation}

\begin{theorem}[Universal mixed-jet floor]\label{thm:main}
Over an algebraically closed field of arbitrary characteristic, every
nonempty finite reduced point-span $s$-osculating-absorbing set for the
complete $H^m$-embedding satisfies
\[
 \dim S_Z\geq B_{d,s}(m),\qquad |Z|\geq B_{d,s}(m).
\]
\end{theorem}

\begin{corollary}[First-order specialization]\label{cor:first}
For $s=1$,
\[
 B_{d,1}(m)=
 \left\lfloor\frac{m+1}{2}\right\rfloor(d+1)
 +\begin{cases}1,&m\text{ even},\\0,&m\text{ odd}.\end{cases}
\]
Thus the formula recovers the growing first-jet bound of
\cite{ErikssonARR2026}.
\end{corollary}

For fixed $d,s$,
\[
 B_{d,s}(m)=\frac{1}{s+1}\binom{d+s}{d}m+O_{d,s}(1).
\]
For $d\geq2$ the leading coefficient strictly increases with $s$, because
\[
 \frac{\binom{d+s+1}{d}/(s+2)}
 {\binom{d+s}{d}/(s+1)}
 =\frac{d+s+1}{s+2}>1.
\]
The hierarchy is therefore asymptotically stronger under the correspondingly
stronger osculating hypothesis.

\section{The annihilator criterion}

Let $(s+1)Z=\coprod_{p\in Z}(s+1)p$.  Finite-scheme duality gives
\[
 S_Z^\perp=H^0(X,\mathcal I_Z\otimes H^m)
\]
and
\[
 \left(\sum_{p\in Z}\widehat\Osc^s_p(H^m)\right)^\perp
 =H^0(X,\mathcal I_{(s+1)Z}\otimes H^m).
\]

\begin{lemma}[Higher-order kernel criterion]\label{lem:kernel}
The set $Z$ is point-span $s$-osculating-absorbing if and only if
\[
 H^0(X,\mathcal I_Z\otimes H^m)
 =H^0(X,\mathcal I_{(s+1)Z}\otimes H^m).
\]
\end{lemma}

\begin{proof}
The inclusion from right to left always holds because
$\mathcal I_{(s+1)Z}\subseteq\mathcal I_Z$.  Absorption is equivalent to
\[
 \sum_{p\in Z}\widehat\Osc^s_p(H^m)\subseteq S_Z.
\]
Taking annihilators reverses this inclusion and supplies the converse
inclusion between the displayed spaces of sections.
\end{proof}

The principal-parts filtration gives
$\widehat\Osc^a_p(H^m)\subseteq\widehat\Osc^s_p(H^m)$ for $a\leq s$.
Thus order-$s$ absorption contains every lower-order jet block.  The converse
is not asserted.

\section{Mixed-jet interpolation}

\begin{lemma}[Full jets from powers]\label{lem:single}
If $p\in X$ and $0\leq a\leq k$, then
\[
 H^0(X,H^k)\longrightarrow
 H^0\!\left((a+1)p,H^k|_{(a+1)p}\right)
\]
is surjective.  Hence
$\dim\widehat\Osc^a_p(H^k)=\lambda_d(a)$.
\end{lemma}

\begin{proof}
Choose $u_0\in H^0(X,H)$ nonzero at $p$.  Very ampleness and smoothness give
$u_1,\ldots,u_d$, vanishing at $p$, such that $x_i=u_i/u_0$ form a regular
system of parameters.  For every multi-index $\alpha$ with $|\alpha|\leq a$,
\[
 u_0^{k-|\alpha|}u_1^{\alpha_1}\cdots u_d^{\alpha_d}
\]
has local representative $x^\alpha$ after trivializing by $u_0^k$.
These monomials form a basis modulo $\mathfrak m_p^{a+1}$.  No
differentiation is used.
\end{proof}

\begin{lemma}[Mixed-jet interpolation]\label{lem:mixed}
Let $t\geq1$, let $p_1,\ldots,p_t$ be distinct points, let $a_i\geq0$, and put
\[
 K=\sum_{i=1}^t(a_i+1)-1.
\]
Then restriction is surjective:
\[
 H^0(X,H^K)\longrightarrow
 \bigoplus_{i=1}^t
 H^0\!\left((a_i+1)p_i,H^K|_{(a_i+1)p_i}\right).
\]
The same holds for $H^k$ in every degree $k\geq K$.
\end{lemma}

\begin{proof}
For each $i\ne j$, choose
$\ell_{ij}\in H^0(X,H)$ with
$\ell_{ij}(p_j)=0$ and $\ell_{ij}(p_i)\ne0$, and set
\[
 P_i=\prod_{j\ne i}\ell_{ij}^{a_j+1}.
\]
It is a unit at $p_i$ and vanishes to the prescribed order at every other
support.  Lemma~\ref{lem:single} realizes arbitrary order-$a_i$ data using
$H^{a_i}$.  Multiplication by the unit $P_i$ is an automorphism of the
truncated local algebra, so degree-$K$ sections realize arbitrary data at
$p_i$ and zero data elsewhere.  Sum over $i$.

For $k>K$, multiply by a section of $H^{k-K}$ nonzero at every support.
It exists because $\kk$ is infinite and a finite union of proper linear
subspaces cannot exhaust the section space.
\end{proof}

Over the complex numbers this is also the line-bundle specialization of
\cite[Proposition~2.3]{BDRS1999}.  The direct proof records the exact
specialization and is characteristic-free.  The interpolation statement itself
is not claimed as new; the point here is its rank consequence under the
absorption constraint and the ensuing optimization.

\section{Convex packing and the rank proof}

The interpolation lemma first gives a heterogeneous certificate, before
optimization.

\begin{proposition}[Mixed-block rank certificate]\label{prop:certificate}
Under the hypotheses of Theorem~\ref{thm:main}, choose $t\geq1$ distinct supports
$p_1,\ldots,p_t\in Z$ and orders $0\leq a_i\leq s$.  If
\[
 \sum_{i=1}^t(a_i+1)\leq m+1,
\]
then
\[
 \dim S_Z\geq\sum_{i=1}^t\lambda_d(a_i).
\]
\end{proposition}

\begin{proof}
Put $K=\sum_i(a_i+1)-1\leq m$.  Lemma~\ref{lem:mixed}, followed by degree
raising when $K<m$, says that the selected infinitesimal neighborhoods impose
independent conditions on $H^m$.  Dually, their osculating blocks form a direct
sum of dimension $\sum_i\lambda_d(a_i)$.  Each block lies in $S_Z$ because
$a_i\leq s$ and $Z$ is order-$s$ absorbing.
\end{proof}

\begin{lemma}[Discrete convex packing]\label{lem:packing}
If $1\leq w_i\leq s+1$ and $\sum_iw_i\leq m+1$, then
\[
 \sum_i\binom{d+w_i-1}{d}\leq B_{d,s}(m).
\]
Equality is attained by $q$ weights $s+1$ and, if $r>0$, one weight $r$.
\end{lemma}

\begin{proof}
Put $f(0)=0$ and $f(w)=\binom{d+w-1}{d}$ for $w\geq1$.  Its forward differences
\[
 f(w+1)-f(w)=\binom{d+w-1}{d-1}
\]
are nondecreasing for $w\geq1$, while $f(1)-f(0)=1$.  Moving one unit from a
smaller block to a larger block below the cap $s+1$ does not decrease the
total; a weight-one block disappears when it loses its unit.  Repetition packs
the budget into full blocks and at most one residual block.  If the original total
weight is smaller than $m+1$, first enlarge a nonfull block or append a
weight-one block; either operation strictly increases the objective.  Thus an
optimum uses the whole budget, and the packed pattern gives \eqref{eq:B}.
\end{proof}

\begin{proof}[Proof of Theorem~\ref{thm:main}]
Put $n=|Z|$.  If $r=0$ and $q=1$, nonemptiness gives $n\geq1=q$.  Now
suppose $r=0$, $q\geq2$, and $n\leq q-1$.  Then $n\geq1$, and the union of the full
order-$s$ neighborhoods at all supports is independent in degree
\[
 (s+1)n-1\leq(s+1)(q-1)-1<m
\]
by Lemma~\ref{lem:mixed}, and remains so after raising degree to $m$.
Its dual has dimension $n\lambda_d(s)$ and lies in $S_Z$ by absorption.
This contradicts $\dim S_Z\leq n$, since $\lambda_d(s)>1$.  Hence $n\geq q$.

If $r>0$ and $n\leq q$, the same contradiction applies because
\[
 (s+1)n-1\leq q(s+1)-1=m-r<m.
\]
Thus $n\geq q+1$.

Choose $q$ supports with full order-$s$ blocks and, if $r>0$, one further
support with an order-$(r-1)$ block.  Their mixed interpolation degree is
\[
 q(s+1)+r-1=m.
\]
Their independent dual spaces lie in $S_Z$ and have total dimension
$B_{d,s}(m)$.  Therefore $\dim S_Z\geq B_{d,s}(m)$, and
$|Z|\geq\dim S_Z$ because $S_Z$ is generated by $|Z|$ value lines.
\end{proof}

The certificate and packing argument do not intrinsically require a tensor
power.  The following form records the more general input precisely.

\begin{corollary}[Jet-ample polarization]\label{cor:jetample}
Let $L$ be an $M$-jet ample line bundle on $X$, meaning that for every set of
distinct points $p_i$ and positive integers $b_i$ with
$\sum_i b_i=M+1$, the restriction map
\[
 H^0(X,L)\longrightarrow
 \bigoplus_i H^0\!\left(b_i p_i,L|_{b_i p_i}\right)
\]
is surjective.  Fix $1\leq s\leq M$, write
$M+1=Q(s+1)+R$ with $0\leq R\leq s$, and define $S_Z$ and the osculating
spaces using $L$ in place of $H^m$.  If a nonempty finite reduced set $Z$ is
point-span $s$-osculating-absorbing, then
\[
 \dim S_Z,\ |Z|\geq
 Q\binom{d+s}{d}+\binom{d+R-1}{d},
\]
where the last summand is interpreted as zero when $R=0$.
Over the complex numbers, repeated application of
\cite[Proposition~2.3]{BDRS1999} gives this conclusion with $M=\ell a$ for
$L=A^{\ell}$ whenever $A$ is $a$-jet ample and $1\leq s\leq\ell a$.
Here $a\geq1$ and $\ell\geq1$ are integers.
\end{corollary}

\begin{proof}
The stated jet-ampleness also gives surjectivity when $\sum_i b_i<M+1$:
add one point outside the selected finite set with the missing positive
weight, and then project away its jet target.  Such a point exists because
$d\geq1$ and the algebraically closed field is infinite.  Thus every selected
collection with $b_i=a_i+1\leq s+1$ and total weight at most $M+1$ is
independent.  The support-threshold proof above applies verbatim with $M$ in
place of $m$, including the separate case $Q=1,R=0$.  Selecting $Q$ full
blocks and, when $R>0$, one residual block gives the displayed rank.  The
last assertion is the tensor-product theorem cited above, iterated $\ell-1$
times.
\end{proof}

\begin{proof}[Proof of Corollary~\ref{cor:first}]
Divide $m+1$ by two.  An odd $m$ has no residual block; an even $m$ has one
residual simple value.  Substitution in \eqref{eq:B} gives the formula.
\end{proof}

\section{Sharpness on the rational normal curve}

\begin{theorem}[Exact curve minimum]\label{thm:curve}
Let $X=\PP^1$, $H=\cO_{\PP^1}(1)$, and $1\leq s\leq m$.  The minimum size
and minimum span dimension of a point-span $s$-osculating-absorbing set for
the complete $m$th Veronese embedding are both $m+1$.
\end{theorem}

\begin{proof}
Since $\lambda_1(a)=a+1$, Theorem~\ref{thm:main} gives
$B_{1,s}(m)=q(s+1)+r=m+1$.

Conversely, choose any $m+1$ distinct points $Z\subset\PP^1$.  A nonzero
binary form of degree $m$ cannot vanish at all of them, so evaluation
\[
 H^0(\PP^1,\cO(m))\longrightarrow H^0(Z,\cO_Z(m))
\]
is an isomorphism.  Hence $S_Z=H^0(\PP^1,\cO(m))^*$, which contains every
osculating space, and $\dim S_Z=|Z|=m+1$.
\end{proof}

There is a second exact regime in every dimension.

\begin{theorem}[Exact top-order minimum]\label{thm:top-order}
Let $X=\PP^d$, $H=\cO_{\PP^d}(1)$, and $s=m\geq1$.  The minimum size and
minimum span dimension of a point-span $m$-osculating-absorbing set are both
\(\binom{d+m}{d}\), over every algebraically closed field.
\end{theorem}

\begin{proof}
Theorem~\ref{thm:main} gives the lower bound because $q=1$, $r=0$, and
\(B_{d,m}(m)=\lambda_d(m)=\binom{d+m}{d}\).  Conversely, the evaluation
lines of all closed points of $\PP^d$ span
\(H^0(\PP^d,\cO(m))^*\): their annihilator would otherwise contain a
nonzero homogeneous form vanishing identically on $\PP^d$.  Select a basis
of evaluation lines and let $Z$ be its supports.  Then
\(|Z|=\dim S_Z=\binom{d+m}{d}\) and $S_Z$ is the whole dual space, so it
contains every order-$m$ osculating space.
\end{proof}

The exact replay uses rational Vandermonde matrices for
$(m,s)=(3,1),(5,2),(8,3),(10,5)$.  It also exhaustively checks 250 finite
convex-packing cases and two simplex-lattice fixtures in the top-order regime.
These computations certify only the displayed finite fixtures and integer
identities, not the universal theorem or priority.

\section{Context and limitations}

Jet ampleness encodes simultaneous interpolation on fat points.  A
tensor-product generation result over the complex numbers appears in
\cite{BDRS1999}.  Higher osculating spaces and fundamental forms are
formulated through principal-parts sheaves over algebraically closed fields
in \cite{MallavibarrenaPiene2024}; that is the intrinsic language used here.

Osculating spaces, higher Gauss maps, and secant varieties have a substantial
literature \cite{BallicoEtAl2004,BallicoFontanari2004,BernardiEtAl2009}.
The Veronese literature in particular relates osculating secant defectivity
to Hilbert functions of fat points.  These works provide context; this
selective comparison does not establish priority for the value-span
absorption floor \eqref{eq:B}.

The authorial note B215 contains the order-two mixed interpolation pattern
\cite{ErikssonB215}.  It is an antecedent, not independent validation.  The
present argument removes its unrelated conditional framework, states the
arbitrary-order absorption hypothesis, works in arbitrary characteristic,
optimizes the closed formula, and proves curve sharpness.

The limitations are material.
\begin{itemize}
\item The higher-order absorption hypothesis is explicit.  No implication
from first-order tangent absorption to second- or higher-order absorption is
asserted or used.
\item The restriction $s\leq m$ ensures the full local rank in
Lemma~\ref{lem:single}.  No formula for $s>m$ is asserted.
\item Completeness of the $H^m$ system is used because all separator products
must be available.  Arbitrary subsystems are outside the theorem.
\item Equality is settled for the numerical minimum in dimension one and for
$s=m$ on projective space.  No equality classification in the remaining
higher-dimensional regimes is claimed; in particular, no nontrivial sharpness
example is supplied here for $d\geq2$ and $s<m$.
\item The literature search is selective, the replay is not a proof checker,
and no independent human peer review has yet been performed.
\end{itemize}

\section{Conclusion}

Once the higher-order hypothesis is stated explicitly, point-span
osculating absorption has a uniform characteristic-free rank consequence.
Very ampleness isolates mixed fat-point blocks of total weight at most
$m+1$, and convex packing gives the explicit floor $B_{d,s}(m)$.  The same
certificate yields a bound for arbitrary jet-ample polarizations.  The
first-order result is recovered exactly, while the full hierarchy is sharp on
rational normal curves and the top-order endpoint is sharp on projective
space in every dimension.  Equality away from these regimes and the geometry
of special higher Gauss fibers remain open.

{\small
\begin{thebibliography}{99}
\bibitem{BallicoEtAl2004}
E.~Ballico, C.~Bocci, E.~Carlini, and C.~Fontanari,
\emph{Osculating spaces to secant varieties},
arXiv:math/0406322 (2004).
\url{https://arxiv.org/abs/math/0406322}

\bibitem{BallicoFontanari2004}
E.~Ballico and C.~Fontanari,
\emph{On the osculatory behaviour of higher dimensional projective
varieties}, arXiv:math/0406319 (2004).
\url{https://arxiv.org/abs/math/0406319}

\bibitem{BDRS1999}
M.~C.~Beltrametti, S.~Di Rocco, and A.~J.~Sommese,
\emph{On generation of jets for vector bundles},
Rev. Mat. Complut. \textbf{12} (1999), 27--45.
\url{https://doi.org/10.5209/rev_REMA.1999.v12.n1.17182}

\bibitem{BernardiEtAl2009}
A.~Bernardi, M.~V.~Catalisano, A.~Gimigliano, and M.~Id\`a,
\emph{Secant varieties to osculating varieties of Veronese embeddings of
$\PP^n$}, J. Algebra \textbf{321} (2009), 982--1004.
\url{https://doi.org/10.1016/j.jalgebra.2008.10.020}

\bibitem{ErikssonARR2026}
L.~Eriksson,
\emph{Characteristic-Free Bounds and Sharp Gauss-Fiber Examples for
Point-Span Tangent Absorption},
ARR-2026-19BXX42GM38K6VZ3, v1 (2026).
\url{https://arr-research.github.io/papers/ARR-2026-19BXX42GM38K6VZ3/}

\bibitem{ErikssonB215}
L.~Eriksson,
\emph{B215---Simultaneous second-jet interpolation raises the floor},
public working note (2026),
\href{https://github.com/lluiseriksson/hodge-conjecture-research-front/blob/main/proofs/B215-simultaneous-second-jet-interpolation-floor.md}{repository record B215}.

\bibitem{MallavibarrenaPiene2024}
R.~Mallavibarrena and R.~Piene,
\emph{On fundamental forms and osculating bundles},
Rend. Istit. Mat. Univ. Trieste \textbf{56} (2024), Art.~9, 21 pp.
\url{https://doi.org/10.13137/2464-8728/36877}
\end{thebibliography}
}

\end{document}
